\summarychapter{19}{Lifted in a Mystic Manner}{The Papacy and the Crusader}

\begin{inbrief}
In the twelfth and thirteenth centuries the papacy tried to build a perfect society on earth and rose
to a power it never held before or since. Like the Gothic cathedrals of the same years it reached for
the impossible, cracked and fell. The chapter shows the papal monarchy under Innocent III and then the
crusades: their causes, their course and their meagre results.
\question{How much may Christians expect of human society, and what would the earth be like if God's will were done there as it is in heaven?}
\end{inbrief}

Notre Dame in Paris was built between 1163 and 1235, while students argued theology and traded news of
the crusades on the nearby bridge; one of them became Pope Innocent III. It was one symptom of a fever
of building: between 1170 and 1270 more than five hundred great Gothic churches rose in France alone.
The style first appeared in the abbey church of St.~Denis, rebuilt between 1137 and 1144 by Abbot
Suger, who felt himself ``lifted in a mystic manner'' from the lower toward the upper sphere. Lift,
reach, aspiration are the words for an age that was far from dark.

\ssection{Thy Kingdom Come}
Gothic builders sought the greatest possible height. Chartres rises as high as a thirty-story
building, Strasbourg as high as forty, and towers often fell. Flying buttresses replaced massive walls
with slender pillars, and even the carved saints were elongated so that everything seemed to soar. So
with the papacy. As unified national states reduced the empire to the sovereignty of a cluster of
German kingdoms and city republics, the papacy, building on Gregory VII's reforms, became the most
powerful office in Europe, a centralized universal monarchy. Every bishop took an oath of loyalty to the pope; founding a religious order required his consent; his court heard appeals from all Christendom; his
legates supervised his orders in every country.

Innocent III (1198--1216) was a new kind of pope, trained not as a monk but as a canon lawyer. The
successor of Peter, he said, is the Vicar of Christ, a mediator ``below God but beyond man,'' who
judges all and is judged by none. The papacy is the sun and kings are the moon, shining with borrowed
light. His weapons were spiritual, and effective because hardly anyone doubted heaven and hell or the pope's control of the grace needed to reach one and escape the other. Excommunication cut
an individual off from the church: a bell tolled, a book was closed, a candle was put out. The
excommunicated could not be judge, witness or party to a contract and were refused Christian burial.
The interdict fell on whole nations, suspending public worship and all sacraments except baptism and
extreme unction; Innocent used or threatened it eighty-five times. Many states acknowledged the pope
as feudal lord. King John of England, at odds with him over the archbishop of Canterbury, was
excommunicated and his realm put under interdict, until he became the pope's vassal, received England
back as a fief and paid tribute. Philip Augustus of France was forced to take back the queen he had
divorced. In Germany Innocent intervened in a civil war and secured the throne for his ward Frederick
II. Under him and his successors Christianity reached its peak of political and cultural influence.
Its soaring dreams, and delusions, appear above all in the crusades and in Scholasticism (the subject
of the next chapter).

\ssection{Taking Up the Cross for Christ}
Islam began with Muhammad, born in 570. A keeper of camels, he came at twenty-five to manage the
affairs of a wealthy widow fifteen years his senior, who married him and was his only wife while she
lived. At forty he began to withdraw to a cave to meditate, and reported that an angel had commanded
him ``to recite''; hence the Koran. His wife was his first follower. After battles and wars he won a
vast following in Arabia, and within a hundred years Jerusalem, Damascus and Cairo were Muslim. His
death in 632 without a named successor divided the movement: Sunnis held that the caliph should be
elected, Shiites that he must come from the prophet's bloodline, beginning with Ali, his cousin and
son-in-law. By the eighth century half of all Christians lived under Islamic rule. Whatever other
motives are suggested---adventure, gain, spiritual advantage---and however few would now defend
the crusades, Shelley insists on one plain fact: Christians were trying to counter a remarkable
military conquest and keep their remaining strongholds.

Pilgrims had long traveled to the Holy Land, and Islam at first did not stop them; a German
pilgrimage of 1065 numbered about seven thousand. In the eleventh century the Seljuk Turks, new and
fanatical converts, seized Jerusalem from fellow Muslims and swept into Asia Minor. At Manzikert in
1071 they captured the eastern emperor and scattered his army, and Byzantium lost its main source of
revenue and troops. Tales of mistreated pilgrims, partly propaganda, inflamed Europe. In 1095, after
an appeal from the emperor Alexius I, Pope Urban II proclaimed the First Crusade at the Council of
Clermont, promising material as well as spiritual gain: your land is too narrow and poor; take the
road to the Holy Sepulchre and win that land for yourselves. The crowd roared ``\emph{Deus vult!}''
(``God wills it!''), which became the battle cry. Every great and colorful figure of the age was caught
up in the cause, from Peter the Hermit, who inflamed the First Crusade, to the saintly Louis IX of
France, who inspired the Sixth and Seventh.

The crusaders' motives were mixed. ``Crusade'' means taking the cross, which they wore on the breast
going out and on the back coming home. They knew the promised rewards, including full forgiveness of
past sins, and most felt deep reverence for the ground Christ had walked. For the popes it was a new
kind of war, a holy war. Augustine's rules for a just war---waged by the state, to defend life and
property, with respect for noncombatants and prisoners---evaporated: the church called the war, its
aim turned from freeing fellow Christians to conquering the infidel, and the cause excused every
cruelty. The crusades began with attacks on Jews; fellow Christians were robbed and raped; atrocities
against Muslims went as far as cannibalism.

\ssection{Events of the Crusades}
Between the late eleventh and the late thirteenth century the popes launched seven major crusades and
many smaller expeditions. The First, made up of nobles from France, Germany and Norman southern Italy,
went overland to Constantinople, where Alexius, who had expected mercenaries, hurried the unruly
horde on. With no more than five thousand men it beat the divided Turks and took Jerusalem amid a
slaughter in which, a chronicler wrote, men rode in blood to their knees; that evening the victors
knelt weeping for joy at the Holy Sepulchre. They founded the Latin kingdom of Jerusalem, whose last
remnant fell in 1291. When the kingdom tottered in 1147, Bernard of Clairvaux preached the Second
Crusade, which achieved nothing and melted away after two years.

The first zeal had cooled, and the cause was being corrupted by money. In medieval practice a
confessed sin was forgiven, but the priest imposed a penance; what was not performed in life was
completed in purgatory. This was ``temporal'' punishment, and the church claimed power to remit part
of it. Urban at Clermont offered the first complete remission, or indulgence, to those who went to
Jerusalem out of pure devotion. It was a small step to grant the same to those who paid for a
substitute, and fund-raising for hospitals and cathedrals followed.

In 1187 Saladin, sultan of Egypt and Syria, retook Jerusalem. The Third Crusade (1189) was led by
three famous kings. Frederick Barbarossa drowned in Asia Minor, Philip Augustus quarreled with
Richard the Lion-Hearted and went home, and Richard was left to face Saladin, who proclaimed a
\emph{jihad} but remained a patient and chivalrous statesman; he even proposed that Richard marry
his sister and receive Palestine as a wedding gift. They agreed on a three-year truce and free access
for pilgrims, which Saladin would have granted anyway. The Fourth Crusade, called by Innocent III,
showed how far the ideal had sunk. Unable to pay Venice for shipping, the knights sacked the Christian
town of Zara for the Venetians in 1202 and were excommunicated; then they took and ravaged
Constantinople itself, set up a Latin Empire there in 1204 and forgot the Holy Land. Innocent raged at
the outrage, then appointed a Latin archbishop. The Latin Empire lasted until 1261; the conquest
widened the schism with the Greek church and hastened the city's fall to the Turks in 1453. Later
crusades changed nothing, and the era ended in 1291 when Acre, the last Christian stronghold, fell.

\ssection{Results of the Holy Cause}
If the aims were to win the Holy Land, check Islam and heal the schism with the East, the crusades
failed miserably. They created three military orders that combined monasticism and warfare to protect
pilgrims and fight Muslims: the Templars (white robe, red cross), the Hospitalers (black robe, white
cross) and the German Teutonic Knights (white robe, black cross), whose castles guarded the roads of
the Crusader States. (The chapter carries two charts, ``The Age of Crusades'' and ``The High
Middle Ages.'') The most significant result was added splendor for the papacy, which launched and inspired the
expeditions and for a time had the new orders and the bishops of the Holy Land and Constantinople
under its protection. In the holy war the popes were grasping at rule over a single church of East and West. But Europe took Syria and Jerusalem in a fit of enthusiasm and could not keep them, and
when trade and towns turned people to new interests the popes clung to the old ideal. They never
grasped, Shelley concludes, two truths: the deepest joys of the faith are not tied to holy sites, and God does not spread Christ's church by force of arms.

\begin{keyfacts}
\keypeople{Innocent III; Abbot Suger; King John of England; Philip Augustus; Frederick II; Muhammad;
Urban II; Alexius I; Peter the Hermit; Bernard of Clairvaux; Saladin; Richard the Lion-Hearted;
Frederick Barbarossa; Louis IX.}
\keydates{%
  \kd{570}{Birth of Muhammad}
  \kd{632}{Death of Muhammad}
  \kd{1071}{Battle of Manzikert: the Seljuk Turks defeat Byzantium}
  \kd{1095}{Urban II proclaims the First Crusade at Clermont}
  \kd{1147}{Second Crusade, preached by Bernard of Clairvaux}
  \kd{1163--1235}{Notre Dame of Paris built}
  \kd{1187}{Saladin recaptures Jerusalem}
  \kd{1189}{Third Crusade}
  \kd{1198--1216}{Pontificate of Innocent III}
  \kd{1204}{Fourth Crusade captures Constantinople; Latin Empire founded}
  \kd{1291}{Fall of Acre: end of the crusades}}
\keyterms{Gothic; Vicar of Christ; excommunication; interdict; canon law; crusade; holy war and just
war; indulgence; temporal punishment; \emph{jihad}; Sunni and Shiite; Templars, Hospitalers, Teutonic
Knights.}
\end{keyfacts}
